\section{Evaluación experimental ISO/IEC 25010}
\label{sec:iso25010}

El modelo de calidad ISO/IEC 25010:2023~\cite{iso25010} organiza la calidad de producto en ocho
características; esta entrega evalúa las cinco exigidas por la guía, cada una con al menos una
métrica cuantitativa y su umbral objetivo (Tabla~\ref{tab:iso-umbrales}), siguiendo el rigor de
reporte experimental de Jain~\cite{jain1991art} y Wohlin et al.~\cite{wohlin2012experimentation}.

\subsection{Preguntas de investigación}
\label{sec:preguntas-investigacion}

Siguiendo el criterio de Wohlin et al.~\cite{wohlin2012experimentation} de que una pregunta de
investigación debe admitir una respuesta negativa para ser una pregunta y no una afirmación
disfrazada, esta entrega formula tres, cada una respondida con datos ya obtenidos y no con una
expectativa:

\begin{enumerate}[nosep]
  \item \textbf{PI1.} ¿El sistema, bajo la carga concurrente medida, cumple los umbrales
        objetivo de eficiencia de desempeño (p95 $<$ 500\,ms) y fiabilidad (error 5xx $<$ 1\,\%)
        de ISO/IEC 25010? \textbf{Respuesta: no}, en ninguna de las dos —ver
        Sección~\ref{sec:iso25010} de este mismo documento (subsección \emph{Eficiencia de
        desempeño y fiabilidad}), con causa raíz identificada y no solo el síntoma.
  \item \textbf{PI2.} ¿El sistema mitiga las seis categorías de riesgo OWASP más aplicables a su
        arquitectura REST + JWT? \textbf{Respuesta: no completamente} —5 de 6, con una brecha
        real y concreta documentada (ausencia de límite de intentos de inicio de sesión), no
        cinco de seis presentadas como éxito total.
  \item \textbf{PI3.} Antes de esta auditoría, ¿el pipeline de CI/CD garantizaba que ningún
        artefacto se publicara en el registro sin que las cinco verificaciones automatizadas
        pasaran primero? \textbf{Respuesta: no} —\texttt{build-images} solo dependía de dos de
        los cinco \emph{jobs} de verificación, y corría en paralelo con el sexto
        (\texttt{integration}) en vez de después. Ver Sección~\ref{sec:pruebas-cicd} para el
        diagnóstico completo y la corrección aplicada.
\end{enumerate}

Las tres preguntas podían haberse respondido \enquote{sí} y el resultado seguiría siendo válido y
reportable: el objetivo no era producir un resultado favorable, sino uno verificable.

\subsection{Nota sobre el alcance}

Por restricción de tiempo de entrega, el equipo decidió correr una \textbf{escala reducida} del
experimento de carga: \textbf{r=5 repeticiones} por escenario (en vez de las 10 exigidas), con
corridas más cortas (60--120\,s en vez de 5--10\,min). La metodología de análisis —descartar la
primera y la última repetición, calcular media e intervalo de confianza al 95\,\% sobre las
intermedias— es la misma que exige el protocolo, aplicada a menos muestras. El detalle completo,
con los datos crudos de las diez corridas reales, está en
\texttt{docs/experimentos/evaluacion\_iso25010.md}.

\begin{table}[H]
\centering
\small
\caption{Las cinco características evaluadas, con su umbral objetivo y el resultado medido.}
\label{tab:iso-umbrales}
\begin{tabular}{@{}p{2.9cm}p{3.1cm}p{4.6cm}p{1.3cm}@{}}
\toprule
\textbf{Característica} & \textbf{Umbral} & \textbf{Medido} & \textbf{¿Cumple?} \\
\midrule
Eficiencia de desempeño & p95 $<$ 500\,ms & 1.4--1.6\,s bajo carga concurrente & No \\
Fiabilidad & Error 5xx $<$ 1\,\% & 19--34\,\% bajo carga concurrente & No \\
Seguridad & 100\,\% (acceso + OWASP) & 30/30 control de acceso; 5/6 categorías OWASP mitigadas & Parcial \\
Mantenibilidad & Cobertura $\geq$70\,\% + CCN $<$10 & 96.6\,\% cobertura; CCN medio 1.8 & Sí \\
Compatibilidad & Pasa suite E2E & 9/9 en 3 navegadores; API 26+ verificado por \emph{lint} & Sí \\
\bottomrule
\end{tabular}
\end{table}

\subsection{Eficiencia de desempeño y fiabilidad}

Se ejecutaron dos escenarios de carga con Locust contra el API Gateway: (A) 30 usuarios en carga
plana durante 60\,s, y (B) una rampa de 0 a 60 usuarios durante 120\,s. En ambos, la tasa de
error observada (19\,\% y 34\,\% respectivamente) superó ampliamente el umbral del 1\,\%. La
causa raíz, confirmada directamente en los registros de \texttt{auth-service} y
\texttt{ticket-service}, es un mensaje explícito del motor de base de datos:

\begin{quote}
\ttfamily\small ERROR: No license installed. The maximum number of concurrently open
transactions has been reached.
\end{quote}

Es decir: bajo concurrencia real, \textbf{CockroachDB Core sin licencia} rechaza nuevas
transacciones al superar su límite de transacciones concurrentes abiertas —una limitación
documentada del propio motor en su edición gratuita, no un defecto de la lógica de reintento de
la aplicación (\texttt{TicketWriter.saveWithRetry}, ya cubierta por pruebas unitarias). Con una
licencia comercial, o migrando a CockroachDB Serverless, este límite no existiría.

\subsection{Seguridad}

Se verificó el control de acceso con 30 peticiones reales contra el API Gateway (10 por caso): sin
encabezado \texttt{Authorization}, con un token de formato inválido, y con un usuario autenticado
de rol incorrecto intentando un endpoint de administración —los tres casos fueron rechazados
correctamente el 100\,\% de las veces (401/401/403).

Para la ausencia de vulnerabilidades OWASP~\cite{owasp2021top10}, un escaneo automatizado
completo (OWASP ZAP) quedó fuera del alcance de tiempo de esta entrega; en su lugar se
verificaron directamente, por código y en vivo, las seis categorías más aplicables a esta
arquitectura REST + JWT. Cinco están mitigadas (control de acceso, criptografía —contraseñas con
\texttt{BCryptPasswordEncoder}, JWT firmado HS256—, inyección —consultas parametrizadas vía JPA,
sin SQL nativo concatenado—, configuración de \texttt{actuator} —solo expone
\texttt{health,info,prometheus}—, y registro/monitoreo). Se encontró una \textbf{brecha real}: no
existe límite de intentos de inicio de sesión (\emph{rate-limiting} o bloqueo de cuenta), por lo
que un atacante puede reintentar contraseñas sin restricción.

\subsection{Mantenibilidad}

Ver Sección~\ref{sec:pruebas-cicd}: cobertura de 96.6\,\% de líneas propias (por encima del
objetivo del 70\,\%, sin contar código generado por protobuf/gRPC) y complejidad ciclomática
media de 1.8 (muy por debajo del umbral de 10).

\subsection{Compatibilidad}

La misma suite Playwright de la Sección~\ref{sec:pruebas-cicd} se ejecutó contra tres motores de
renderizado (Chromium, Firefox y WebKit —este último equivalente a Safari, que no corre de forma
nativa en Windows/Linux). Las nueve pruebas pasan en los tres motores. Para la aplicación móvil,
la evidencia disponible en este ciclo es el \texttt{minSdk=26} declarado en Gradle (una garantía
de tiempo de compilación) y el resultado limpio de la regla \texttt{NewApi} de Android
Lint, que falla si el código usa una API de plataforma por encima del mínimo sin una guarda de
versión explícita.

\subsection{Amenazas a la validez de este experimento}

\textbf{Internas}: (1) la escala reducida (5 repeticiones, no 10) amplía los intervalos de
confianza frente al protocolo completo; (2) Locust, el \emph{stack} completo de 17 contenedores y
el propio proceso de verificación compitieron por CPU en una sola máquina de desarrollo,
introduciendo ruido en la latencia medida; (3) el límite de licencia de CockroachDB es una
propiedad de este entorno de demostración, no necesariamente de un despliegue con licencia.
\textbf{Externas}: (1) los datos usados son sintéticos, no tráfico real de un ISP; (2) dos de los
diez tipos de prueba de la pirámide completa (las instrumentadas de Android) no se ejecutaron en
este ciclo por falta de un emulador conectado en el momento de la medición.
